\documentclass[12pt]{article}

\usepackage{amsmath}
\usepackage{graphicx}
\usepackage{natbib}
\usepackage{german}
\DeclareMathOperator\sgn{sgn}

%\graphicspath{{/home/ritz/Teaching/fcn/fig/}}

%----------------- layout ----------------------------------------

\oddsidemargin=0in
\textwidth=6in
\topmargin=-0.7in
\textheight=9.5in
\baselineskip=16pt
\font\mittel=cmr17
\font\roman=cmr12
\font\klein=cmr9
\font\mgross=cmr10
\newfont{\CAP}{cmcsc10 scaled 900}
\pagestyle{empty}
\parindent=0mm

%----------------- macros ----------------------------------------
\newcommand{\conc}[1]{[{\rm C}]_{\text{#1}}}
\newcommand{\units}[1]{\text{#1}}
\newcommand{\K}{\text{K}}
\newcommand{\Na}{\text{Na}}
\newcommand{\Cl}{\text{Cl}}
\newcommand{\Leak}{\text{leak}}
\newcommand{\water}{\text{H}_2\text{O}}
\newcommand{\tderiv}[1]{\frac{\mathrm{d}#1}{\mathrm{d}t}}
%-----------------------------------------------------------------



\renewcommand{\refname}{Literature}
\begin{document}

\parbox{2cm}{
\includegraphics[width=1.8cm]{../bccnlogo.pdf}
}
\parbox{11cm}{
\begin{center}
\large HUMBOLDT-UNIVERSIT"AT \hskip 0.1 cm ZU \hskip 0.1 cm BERLIN
\vskip 0.1 true cm
\mgross BERNSTEIN CENTER FOR\\ COMPUTATIONAL NEUROSCIENCE
\end{center}
}
\parbox{2cm}
{
\hfill
\includegraphics[width=1.8cm]{../hublogo.pdf}
}

\vskip 0.6 true cm

\leftline{\CAP Humboldt-Universit"at zu Berlin 
\hfill Phone: 030/2093-9110}
\leftline{\CAP Philippstr. 13 House 6
              \hfill Fax: 030/2093-6771}
\leftline{\CAP \hfill webpage: http://www.bccn-berlin.de/}


\vskip 0.8 true cm
\centerline{\bf Models of Neural Systems, WS 2010/11}
\centerline{\bf Project 9: Temporal Pattern Learning}
\centerline{Project presentation and report submission: February,
15th, 2011}

\vskip 0.8 true cm

{\bf Background}

The tempotron, in analogy the perceptron, is a neuron that is able to
distinguish between temporal input patterns. It should
respond with an action potential for some input patterns, but not for
others. Importantly, for all input patterns, the tempotron can receive
input from the same cells, but only the relative timing of these inputs is
different for different input patterns. Read the original paper
\citep{Guetig06} to understand the details on the
tempotron algorithm.

\vskip 0.8 true cm

{\bf The model neuron and its inputs} 

The tempotron is a leaky integrate-and-fire neuron that receives input
from $n$ cells. An input spike evokes a postsynaptic potential in the
membrane potential of the tempotron. The summed postsynaptic
potentials of all inputs define the membrane potential of the
tempotron: $$V(t)=\sum_i \omega_i \sum_{t_i}\kappa (t - t_i) +
V_{rest}$$ with $t_i$ being the spike time of the $i$th input cell and
$\kappa (t - t_i)$ being the normalized postsynaptic potential evoked
by an input spike (for $t-t_i \ge0$, otherwise $\kappa =0$):

$$\kappa (t -
t_i)=V_0\cdot\left[e^{-(t-t_i)/\tau}-e^{-(t-t_i)/\tau_s}\right],$$
where $V_0$ is a normalization constant to ensure that $\kappa (t -
t_i)$ integrates to 1; $\tau$ and $\tau_s$ are time constants for the membrane
and synaptic currents,  which determine the decay and
rise time of the post-synaptic potentials. For $\tau/\tau_s = 4$, the
normalization constant is  $V_0=2.12$.

\begin{enumerate}
   \item Implement a basic version of the tempotron. Start with very
   few input cells with different synaptic strengths $\omega_i$. Check
   a few voltage traces to ensure that your implementation is correct.
   \item Generate a set of random input patterns. Each pattern occurs
   within a time window of 500~ms. Each input cell fires only once
   per input pattern. The time of the input spike for each input cell
   is uniformly-distributed within the 500~ms window.
   \item Vary the synaptic weights and the number of input cells.
\end{enumerate}    

{\bf Tempotron learning} 

Now, we want to use a rule that modifies
the synaptic weights when a temporal pattern is presented. The goal is
that the tempotron produces an output spike for some patterns (`+'
patterns) but not for others (`-' patterns). To determine whether an
output spike is produced, we simply check whether the membrane
potential $V(t)$ exceeds a certain threshold for a given pattern.
Intuitively, the synaptic weights should be increased if we want our
tempotron to fire an action potential for a certain pattern.
Conversely, the synaptic weights should be decreased if we do not want
to fire a spike for a certain pattern. However, as the input patterns
mainly differ in the timing of the input, we need a sophisticated
learning rule. The tempotron learning rule is
  $$ \Delta\omega_i = \lambda \sum_{t_i<t_{max}} \kappa
  (t_{max}-t_i)$$ where $\lambda>0$ is a learning rate and $t_{max}$
  is the time when $V(t)$ is maximal. Weights are only changed for
  incorrect responses of the tempotron. Weights are increased by
  $\Delta\omega_i$ if no spike was elicited (although there should have
  been one). Weights are decreased by $\Delta\omega_i$ if a spike was
  elicited (although there should have been none).

\begin{enumerate}
  \item Train your tempotron on two patterns (one `+' and one `-'
  pattern). Compare the synaptic weights before and after learning.
  Switch the identities of the `+' and `-' pattern and learn again.
  
  \item Change some aspects of the learning paradigm (e.g.~ number of
      patterns, the
  temporal structure of the input patterns, the shape of postsynaptic
  potentials, the learning parameters,...) and examine the effect on
  learning performance \citep[see Figures 3~--~6 in][]{Guetig06}.
  \item For what types of learning could the tempotron learning rule
  be used by the brain? Try to illustrate your idea(s) with
  simulations using the tempotron. Are there problems or weaknesses of
  the algorithm that need to be addressed by real-world neural systems? 

\end{enumerate}

\begin{thebibliography}{9}
    \bibitem[G\"utig and Sompolinsky(2006)]{Guetig06} 
        R.~G\"utig and H.~Sompolinsky, \emph{The tempotron: a neuron that
        learns spike timing-based decisions.} Nature Neuroscience}
        9(3):420--8, 2006
    \bibitem[G\"utig and Sompolinsky(2009)]{Guetig09} 
        R. G\"utig and H. Sompolinsky, \emph{Time-warp-invariant neuronal
        processing}. PLoS Biology 7(7): e1000141, 2009
\end{thebibliography}



\vfill
\centerline{\CAP Contact} \CAP

\begin{tabular}{lll}
Urs Bergmann & Phone: 2093-8924 & Email:
urs.bergmann(at)biologie.hu-berlin.de \\
Matthias Guggenmos & & Email: matthias.guggenmos(at)bccn-berlin.de \\
Richard Kempter \hfill & Phone: 2093-8925 \hfill & Email:
r.kempter(at)biologie.hu-berlin.de \\
Bartosz Telenczuk & Phone: 2093-8838 &
Email: b.telenczuk@biologie.hu-berlin.de \\
\end{tabular}


\end{document}
 
\bibliographystyle{plainnat}
